\documentclass[11pt,a4paper]{article}
\usepackage{amsmath,amssymb,amsthm}
\usepackage[margin=2.5cm]{geometry}
\usepackage{hyperref}

\newtheorem{definition}{Definition}[section]
\newtheorem{theorem}[definition]{Theorem}
\newtheorem{proposition}[definition]{Proposition}
\newtheorem{lemma}[definition]{Lemma}
\newtheorem{corollary}[definition]{Corollary}
\newtheorem{conjecture}[definition]{Conjecture}
\newtheorem{remark}[definition]{Remark}

\title{Hyperseed Formalization:\\
  OpenBGI is Building its Initial Network}
\author{ProtomegaTron (automated formalization)}
\date{2026-09-18}

\begin{document}
\maketitle

\begin{abstract}
We formalize the intellectual structure of Ben Goertzel's Substack article
``OpenBGI is Building its Initial Network'' (2026-05-07) within the
Hyperseed ontology. The article announces a network for deploying
heterodox AGI in a post-LLM-scaling window. We model the cracking of
transformer orthodoxy as curvature revelation in the intelligence
landscape, the multi-paradigm requirement as fiber decomposition with
synergy constraints, and the network organizational form as a
base-space topology resistant to institutional capture.
\end{abstract}

\tableofcontents

%% =================================================================
\section{Scaling curves and curvature revelation}
\label{sec:scaling}
%% =================================================================

\begin{definition}[Intelligence landscape --- claims 1--5]
\label{def:landscape}
Let $\mathcal{L} : \mathbb{R}_+ \to \mathbb{R}_+$ be the function mapping
compute investment $c$ to capability $\mathcal{L}(c)$. The
\emph{transformer-scaling hypothesis} assumed:
\[
  \mathcal{L}(c) \approx \alpha \cdot c^\beta \quad (\beta \approx 1)
\]
i.e., a nearly flat (zero-curvature) landscape where more compute
yields proportional capability gains.
\end{definition}

\begin{proposition}[Empirical curvature revelation --- claim 1]
\label{prop:flattening}
Observed data shows $\beta$ declining across model generations:
\[
  \beta_{n+1} < \beta_n \quad \Rightarrow \quad
  \frac{d^2 \mathcal{L}}{dc^2} < 0
\]
The landscape has \emph{negative curvature} in the scaling direction:
diminishing returns are an intrinsic geometric property of the
transformer architecture in capability space, not a temporary
engineering bottleneck.
\end{proposition}

\begin{corollary}[Post-scaling window --- claim 6]
The curvature revelation creates a strategic window: the field
recognizes that the ``flat landscape'' assumption was wrong, but has
not yet converged on which alternative direction to explore. This
is the 12--18 month window in which heterodox approaches can gain
traction before a new consensus forms.
\end{corollary}

%% =================================================================
\section{Multi-paradigm fiber decomposition}
\label{sec:multi-paradigm}
%% =================================================================

\begin{definition}[Cognitive architecture bundle --- claims 12--15]
\label{def:cog-bundle}
Let $\pi : E \to B$ be a fiber bundle where:
\begin{itemize}
\item $B$ = the space of cognitive tasks (perception, reasoning,
  planning, learning, creativity, \ldots),
\item $F = F_{\text{neural}} \times F_{\text{symbolic}} \times
  F_{\text{evolutionary}}$ = the paradigm fiber,
\item Each task $b \in B$ requires a particular mix of paradigm
  activations from $F$.
\end{itemize}
\end{definition}

\begin{proposition}[Single-paradigm insufficiency --- claim 5]
\label{prop:single-insufficient}
The transformer architecture defines a section
$\sigma_{\text{LLM}} : B \to E$ that projects onto $F_{\text{neural}}$
alone:
\[
  \sigma_{\text{LLM}}(b) = (f_{\text{neural}}(b), 0, 0)
\]
This section is well-defined for tasks $b$ where neural pattern
matching suffices but has zero activation in the symbolic and
evolutionary fibers, making it degenerate for tasks requiring
logical reasoning, sparse learning, or open-ended exploration.
\end{proposition}

\begin{theorem}[Cross-paradigm synergy requirement --- claim 12]
\label{thm:synergy}
For the full task space $B$, no single-paradigm section achieves
non-degenerate coverage. A viable AGI section must activate multiple
fibers with proper transition functions:
\[
  \sigma_{\text{AGI}}(b) = (f_{\text{neural}}(b),\;
    f_{\text{symbolic}}(b),\; f_{\text{evolutionary}}(b))
  \quad \text{with all components } > 0
\]
and the transition functions $\tau_{ij}$ between paradigm fibers
must satisfy the cocycle condition (cognitive synergy):
\[
  \tau_{ij} \circ \tau_{jk} = \tau_{ik}
\]
so that information flows consistently between paradigms.
\end{theorem}

\begin{remark}[Hyperon as concrete section --- claim 15]
The Hyperon project (MeTTa, MORK, PeTTa, MM2) is a concrete
implementation of $\sigma_{\text{AGI}}$: MeTTa provides the symbolic
fiber, neural networks provide the neural fiber, evolutionary
algorithms provide the evolutionary fiber, and the MeTTa interpreter
serves as the transition-function infrastructure that enables
synergistic interaction.
\end{remark}

%% =================================================================
\section{Network topology vs.\ lab topology}
\label{sec:topology}
%% =================================================================

\begin{definition}[Organizational base space --- claims 9--11]
\label{def:org-space}
Define the organizational base space $B_{\text{org}}$ over which
AGI development occurs. Two topologies:
\begin{itemize}
\item $B_{\text{lab}} \cong \{b_1, \ldots, b_k\}$ (finitely many
  isolated labs, $k \approx 3$--$5$): discrete, capturable.
\item $B_{\text{network}}$: a connected graph of diverse participants
  (researchers, operators, industry, capital) with high connectivity
  and no bottleneck nodes.
\end{itemize}
\end{definition}

\begin{proposition}[Network resilience --- claim 9]
\label{prop:resilience}
The network topology $B_{\text{network}}$ is structurally resistant to:
\begin{enumerate}
\item \textbf{Entryism}: no single node controls enough of the network
  to redirect it (cf.\ the EA entryism analysis).
\item \textbf{Regulatory capture}: compliance requirements cannot be
  imposed through a single chokepoint.
\item \textbf{Paradigm lock-in}: diverse participants maintain
  multiple paradigm fibers simultaneously, preventing collapse to
  a single approach.
\end{enumerate}
\end{proposition}

%% =================================================================
\section{Singularity dynamics}
\label{sec:singularity}
%% =================================================================

\begin{definition}[Acceleration measure --- claims 16--18]
\label{def:acceleration}
Let $R(t)$ be the rate of singularity-relevant innovation at time $t$.
The pre-singularity hypothesis states:
\[
  \frac{dR}{dt} > 0 \quad \text{and} \quad
  \frac{d^2 R}{dt^2} > 0
\]
Both the rate and the acceleration of the rate are positive.
\end{definition}

\begin{proposition}[Selective acceleration --- claim 18]
\label{prop:selective}
The acceleration is selective: $R(t)$ increases only for
singularity-relevant technologies, not for all technologies:
\[
  R_{\text{AGI}}(t) \gg R_{\text{socks}}(t) \approx \text{const}
\]
This selectivity is consistent with a phase transition concentrated
in a specific region of technology space, not a general speedup.
\end{proposition}

%% =================================================================
\section{d-calculus perspective}
\label{sec:d-calc}
%% =================================================================

\begin{remark}[Post-scaling as curvature revelation]
In the d-calculus framework, the transformer-scaling era was a period
where the field followed a semantic highway (``just scale more'') that
appeared flat. The empirical flattening reveals that this highway
crosses a region of high curvature: the returns diminish because
the underlying intelligence landscape is not flat in the scaling
direction. The post-LLM window is the moment when the field
collectively realizes the highway is curved and begins exploring
alternative directions.
\end{remark}

\begin{remark}[Network as distributed fiber transport]
The OpenBGI network structure implements distributed fiber transport:
each participant maintains and develops their own paradigm fibers,
and the network provides the communication infrastructure for
synergistic interaction. This is cognitive synergy at the
organizational level, mirroring the cognitive synergy required
at the architectural level within a single AGI system.
\end{remark}

\end{document}
